% Transcription — fonds Alexandre Grothendieck, shelfmark 12, batch 8 (pages 141-160).
% Established from the facsimile archives/batches/12/batch-08.pdf (a mirror of
% https://grothendieck.umontpellier.fr/12.pdf), per the skill
% transcribe-grothendieck. The batch file carries no cover sheet: PDF page N is
% page 140+N, as the Montpellier footer printed on each PDF page confirms, and
% the archivists' pencil numbers bottom-left agree with that position on every
% leaf where they are legible (151 to 160 read plainly).
%
% Pass: Opus 5.5 (claude-opus-5-5), 2026-09-23 — first pass, unchecked against the pages by a human.
%
% What the batch is. All in his hand except page 159. No leaf is dated and he
% does not paginate this run.
%   141     The end of a discussion of smooth motives over a base S: a functor
%           (read \mathcal{L}) is fully faithful; its essential image; whether
%           unramified T_ℓ for all ℓ forces good reduction of an abelian scheme
%           (a criterion of Ogg-Šafarevič type in dim 1, doubtful in mixed
%           characteristic dim 2), and a « Pb » on pulling back smooth motives.
%   142     A cover: otherwise blank, bearing his title « Enveloppes
%           pseudo-abéliennes de catégories multiplicatives (Constructions
%           abstraites de catégories de motifs) ». No \page; the title is the
%           \section.
%   143-147 « Construction abstraite d'une théorie de motifs semi-simples »:
%           an additive Q-linear (k-linear) category V with Fl(X,X) finite-
%           dimensional semisimple; the universal additive functor into an
%           abelian category; the pseudo-abelian envelope M^# of pairs (X, π),
%           π² = π, with the universal property Fonc add(M,C) ≃ Fonc add(V,C)
%           for C with split idempotents; the claim that M is abelian
%           semisimple (idempotent calculus u, v, π, ϖ on p. 146; chain
%           conditions from finite dimension); reduction of (ii) to P = Q.
%   148     Biadditive functors V × V' → V'' extend to M × M' → M''; for a
%           symmetric associative unital ⊠ on V, ⊗ on M, and in char. 0 the
%           tensor operations ⊗^n, ∧^n, S^n.
%   149-151 When the exact T : M → C induced by H is faithful; Example: V from
%           smooth projective varieties over k algebraically closed with
%           Fl(X,Y) = ⊕ C^{m_j}(X_i × Y_j) (cycles mod ℓ-adic homological
%           equivalence); assuming hom = num and the Hodge inequality,
%           End is a Castelnuovo-Weil involutive algebra, hence semisimple;
%           the Tate functor T_ℓ : M → Mod_tf(Q_ℓ); Hodge, De Rham, Betti
%           functors; Künneth ⇒ T_ℓ is a ⊗-functor. Then (p. 151, after a
%           triple rule) back to the abstract situation with a grading by
%           projectors π_i.
%   152-153 π_i ∈ End id_V, Σπ_i = 1 (« somme infinie ! »): M graded by
%           weight, M^i × M^j → M^{i+j}, K(M) a graded λ-ring; a Tate object
%           Λ(1) of degree p ≥ 1 with M ↦ M(1) fully faithful; the category
%           \widetilde{M} of virtual motives M(i), i ∈ Z, as a pseudo-lim of
%           copies of M, periodic of period p.
%   154     « Niveau d'un motif » (his underlined title, over a struck
%           « Décomposition canonique »): level ≤ i − kp, the thick
%           subcategories, M^i ≃ P^i × P^{i-p}(1) × P^{i-2p}(2) × …, with
%           P^j the primitive effective motives, and the induced products.
%   155-156 « Duals et Homs de motifs » (his title): internal Hom, M* =
%           Hom(M, Λ(−i)), Hom(M, N(−i)) ≃ M* ⊗ N, \check{M} = Hom(M, Λ(0));
%           bilinear forms B(M,N), alternating/symmetric, non-degeneracy.
%   157     « Involutions » (his title): the adjoint u' of u : M → N for
%           non-degenerate forms α, β, its rules, End(M) as an involutive
%           Q-algebra, and change of forms u'' = A^{-1} u' B.
%   158     « Formes définies positives sur les motifs » (his title): End
%           h^i(X) are Riemann algebras; sign indeterminacy ε_i of the system
%           φ_X, fixed by polarisations of abelian varieties.
%   159     Unrelated English typescript (a page « VII.19 » on Brauer groups
%           of regular domains), upside down and unmarked by him. Absent.
%   160     The axioms (i)-(iii) on the sets Φ_X of fundamental forms which
%           p. 158 already uses, (iii) being positivity Tr(s_u u) ≥ 0: the
%           leaf is logically prior to 158. Batch 9 opens on its sequel.
%
% Notation fixed for the batch. The categories are \mathcal{V}, \mathcal{M},
% \mathcal{C}, \mathcal{P} (primitive motives, a doubtful reading of a round
% letter); his opposite-category mark is read ^{\circ}; the universal functor
% is \underline{h} (his underline); the internal Hom, written with a heavy
% cursive H, is \underline{\operatorname{Hom}}; the completed category is
% \widehat{\mathcal{M}} on pp. 155-156 (circumflex on the page) and
% \overline{\mathcal{M}} on p. 158 (bar on the page) — kept as written; the
% sets of fundamental forms are \Phi_X, as batch 9 has them; ⊠ is the product
% on V, ⊗ on M.
%
% Legibility. Pages 143-145, 150-155 and 157 are written fair and read nearly
% through; 141, 146, 147, 149, 158 and the lower half of 160 are fast, struck
% over, and their prose is often lost while the formulas carry.
\documentclass[11pt,a4paper]{article}
% Le préambule commun à toutes les transcriptions du fonds Grothendieck.
%
% Il tient en une idée : **l'incertitude se compose, elle ne se résout pas.**
% Une transcription qui lisse un mot douteux a détruit la seule chose pour
% laquelle on la faisait — un lecteur doit pouvoir savoir, à chaque endroit,
% ce qui était sur le feuillet et ce qui a été ajouté. Les six macros qui
% suivent sont donc l'appareil critique, et non de la décoration.
%
% Les mêmes commandes sont reconnues par `scripts/render.mjs`, qui produit la
% vue de lecture du site. Une macro ajoutée ici doit l'être aussi là-bas,
% faute de quoi le rendu échouera bruyamment — ce qui est voulu : un rendu
% silencieusement faux est pire qu'un rendu absent.

\ProvidesPackage{grothendieck}[2026/08/08 Transcriptions du fonds Grothendieck]

\usepackage{amsmath,amssymb,amsthm}
\usepackage{mathrsfs}
\usepackage{stmaryrd}
% Les diagrammes commutatifs sont partout dans ce fonds : c'est la langue
% même de la géométrie algébrique de Grothendieck, et une transcription qui
% les rendrait en prose ne transcrirait rien.
\usepackage{tikz-cd}
% Français par défaut — c'est la langue des feuillets — mais l'anglais est
% chargé aussi : la traduction et le résumé sont des documents anglais, et la
% typographie française (espaces avant les deux-points) y serait une faute.
% Ces documents basculent par \selectlanguage{english} en tête de corps.
\usepackage[english,french]{babel}
\usepackage{fontspec}
\usepackage{geometry}
\geometry{a4paper,margin=25mm,marginparwidth=28mm}
\usepackage{marginnote}
\usepackage{xcolor}
% ulem, not soul. `soul`'s \st and \ul are built on a character-by-character
% scan that XeLaTeX's font machinery breaks: the document fails with an
% undefined control sequence rather than degrading. ulem does the same two jobs
% and is Unicode-engine safe. `normalem` stops it hijacking \emph, which the
% transcriptions use for Grothendieck's own emphasis.
\usepackage[normalem]{ulem}
\usepackage{hyperref}
\hypersetup{colorlinks=true,linkcolor=black,urlcolor=black,pdfborder={0 0 0}}

% --- Métadonnées du lot -----------------------------------------------------
% Reprises telles quelles par le script de rendu, qui les affiche en tête de
% la vue de lecture. Elles fixent ce que le fichier prétend couvrir : sans
% elles, une transcription flotte sans point d'attache dans un fonds de seize
% mille feuillets.

\newcommand{\folder}[1]{\gdef\grothFolder{#1}}
\newcommand{\batch}[1]{\gdef\grothBatch{#1}}
\newcommand{\leaves}[2]{\gdef\grothFirst{#1}\gdef\grothLast{#2}}
\newcommand{\foldertitle}[1]{\gdef\grothFoldertitle{#1}}
\gdef\grothFolder{?}\gdef\grothBatch{?}\gdef\grothFirst{?}\gdef\grothLast{?}\gdef\grothFoldertitle{}

% --- Le résumé --------------------------------------------------------------
%
% Il ouvre la lecture modernisée, sous le titre « Résumé », et il est composé
% en petit. Ce n'est pas un ornement : c'est le seul passage du document qui
% ne soit pas des mathématiques, et un lecteur doit pouvoir le reconnaître
% avant de l'avoir lu — puis le sauter s'il connaît déjà le sujet, ou s'y
% arrêter s'il ne le connaît pas. Le filet qui le suit marque la reprise du
% corps.

\newenvironment{resume}
  {\small\setlength{\parskip}{0.45em}}
  {\par\medskip\hrule\medskip}

% --- Mots-clés ---------------------------------------------------------------
%
% En anglais, en clôture du résumé de la lecture modernisée : le vocabulaire
% moderne sous lequel chercher ce que les feuillets construisent. Le manifeste
% du site (`npm run manifest`) les extrait pour étiqueter la cote entière —
% cette ligne est la source unique des tags, il n'y a pas de fichier à côté.
\newcommand{\keywords}[1]{\par\smallskip\noindent
  {\footnotesize\textsc{Keywords} --- \textit{#1}}\par}

% --- L'appareil critique ----------------------------------------------------

% Le numéro de feuillet, en marge. C'est l'ancre qui permet de régler un
% désaccord sur une formule en regardant *un* feuillet plutôt qu'une chemise
% de six cents. Le site s'en sert aussi pour tourner les pages du fac-similé
% quand on fait défiler la transcription.
\newcommand{\leaf}[1]{%
  \marginnote{\footnotesize\color{gray!70}\textsf{#1}}%
  \label{leaf:#1}\ignorespaces}

% Illisible : on ne devine pas. Un mot inventé qui se lit comme les autres est
% le pire résultat possible.
\newcommand{\ill}{\textcolor{red!60!black}{[\,\ldots\,]}}

% Lecture douteuse : le mot est donné, et signalé comme incertain.
\newcommand{\uncertain}[1]{\uline{#1}}

% Ajout de l'éditeur — une lettre manquante, un mot sous-entendu.
\newcommand{\add}[1]{\textcolor{blue!50!black}{[#1]}}

% Biffé par Grothendieck lui-même. Ce qu'il a rayé fait partie du document :
% une rature dit souvent où la pensée a changé de direction.
\newcommand{\struck}[1]{\sout{#1}}

% Note du transcripteur, sur l'état matériel du feuillet.
\newcommand{\note}[1]{\par\smallskip{\small\color{gray!60!black}\textit{[#1]}}\par\smallskip}

% Note marginale de Grothendieck, distincte du corps du texte.
\newcommand{\marginal}[1]{\par\smallskip\noindent\hspace{1em}%
  {\small\color{gray!60!black}#1}\par\smallskip}

% --- Titre ------------------------------------------------------------------

\AtBeginDocument{%
  \begin{center}
    {\large\bfseries Fonds Alexandre Grothendieck --- cote n\textsuperscript{o}~\grothFolder}\\[2pt]
    {\itshape\grothFoldertitle}\\[4pt]
    {\small feuillets \grothFirst--\grothLast\ (lot \grothBatch)}\\[2pt]
    {\footnotesize\color{gray!60!black}Transcription établie d'après le fac-similé numérisé
     par l'Université de Montpellier.\\
     Les lectures incertaines sont soulignées, les illisibles notées [\,\ldots\,],
     les ajouts entre crochets.}
  \end{center}
  \vspace{1em}}

\endinput


\folder{12}
\batch{8}
\pages{141}{160}
\dating{1964-[vers 1971]}
\watermark{Édition de démonstration}
\foldertitle{Généralités (sous-groupes de Galois motiviques) (1964 ?) : notes manuscrites (s.d.), tapuscrit (s.d.)}

\begin{document}

\page{141}
\textbf{Le foncteur $\mathcal{L}$ est pleinement fidèle} (on se ramène
\uncertain{au cas de la base} $S = \operatorname{Spec} K$).
\note{la lettre cursive nommant le foncteur, lue $\mathcal{L}$, est douteuse ;
elle désigne sans doute le foncteur introduit au lot précédent}

L'image essentielle est \struck{\ill{}} \struck{dans} formée des motifs
$(M, E_{\xi}, \mu_{\xi})$ tels que \struck{les} $E_{t}$ \uncertain{leur}
\uncertain{génération} \ill{} \uncertain{fournissent} \uncertain{soit}
associé à un \emph{schéma abélien} sur $S$. Il faut voir aussi que $M$ soit
effectif de poids 1, i.e.\ qu'il soit défini par un schéma abélien $A_{K}$
sur $K$. \struck{\ill{}} \struck{\ill{}}
\note{« $A_{K}$ » est écrit en interligne au-dessus de « sur $K$ »}

La question si cette condition est suffisante revient à celle-ci : Soit $A$ un
schéma abélien sur $K$, \struck{de} tel que $\forall \ell$, \struck{les}
$T_{\ell}(A)$ soit non ramifié sur $S_{\ell}$, \uncertain{dans} $S$.
\uncertain{Est-ce} que si $A$ a bonne réduction
\ill{} (critère de \uncertain{Ogg-Šafarevič}), \ill{} $S$ est de
dim 1, ou si $S$ de car.\ nulle (critère de \ill{}), mais peut-être pas en
général [p.ex.\ \ill{} \uncertain{inégales} \uncertain{caractéristiques}],
\uncertain{pour} $S$ local régulier de dim 2 d'inégales caractéristiques].
\note{l'interligne « (critère de Ogg-Šafarevič) » est d'une lecture très
douteuse : on devine « O?SS » puis « Chafarevich » ; il est rattaché par un
trait au « dim 1 » de la ligne suivante}
Pour faire un contre-exemple, il suffirait d'avoir un exemple d'un schéma
abélien $A$ sur $S$, et d'un sous-groupe fini de $A_{K}$ qui ne se prolonge pas
en un sous-groupe \emph{fini et plat} de $A$. S'il y a un tel exemple, il y en
aurait aussi avec un ordre premier …

\textbf{Pb} Avec la définition adoptée de motif lisse (en termes de celle de
isomotif lisse, comme un isomotif sur $K$ « galoisien » \uncertain{satisfaisant
simplement une} propriété de non-ramification pour les $T_{\ell}$) il ne semble
pas clair comment définir l'image inverse d'un motif (resp.\ d'un isomotif)
par $S' \to S$, par ex.\ les fibres \ill{} sur $S$. Cette difficulté
\struck{\ill{}} \uncertain{concerne} \ill{} les isomotifs, et vaut déjà pour
ceux effectifs de poids 1.
\note{un double trait vertical de sa main, dans la marge gauche, borde tout ce
paragraphe}

\section{Enveloppes pseudo-abéliennes de catégories multiplicatives
(Constructions abstraites de catégories de motifs)}
\note{titre inscrit de sa main en haut à droite d'un feuillet par ailleurs
vierge, la page 142, qui sert de chemise à la suite}

\page{143}
\textbf{Construction abstraite d'une théorie de motifs semi-simples.}

On part d'une catégorie $\mathcal{V}$ \struck{\ill{}} additive${}^{*}$, en
particulier \struck{\ill{}} les \struck{\ill{}} $\operatorname{Fl}(X,Y)$ sont
munis d'une structure de groupe abélien, les accouplements de composition
étant bilinéaires.
\struck{[On se \ill{} pour que \ill{} \ill{} \ill{} \ill{} \ill{} \ill{} existe
un objet nul, \ill{} \ill{} \ill{} deux objets …]}
\note{un passage entre crochets, encadré et biffé de plusieurs traits obliques ;
on y lit « existe un objet nul » et « deux objets »}
On s'intéresse, pour toute catégorie \emph{abélienne} $\mathcal{C}$, aux
foncteurs \struck{\ill{}} « additifs » (i.e.\ commutant aux sommes finies${}^{2}$,
\uncertain{ou encore} \uncertain{qui commutent aux produits finis})
\[
  H \colon \mathcal{V}^{\circ} \longrightarrow \mathcal{C}
\]
et on cherche un foncteur universel de cette nature
\[
  \underline{h} \colon \mathcal{V}^{\circ} \longrightarrow
  \mathcal{M}(\mathcal{V}) = \mathcal{M}
\]
($\underline{h}$ foncteur additif dans une catégorie abélienne $\mathcal{M}$) tel
que pour toute catégorie abélienne $\mathcal{C}$, l'application
\[
  T \longmapsto T \circ \underline{h} \colon\quad
  \operatorname{Fonc\,exacts}(\mathcal{M}, \mathcal{C}) \longrightarrow
  \operatorname{Fonc\,add}(\mathcal{V}^{\circ}, \mathcal{C})
\]
soit bijective (ce qui caractérise $\mathcal{M}$ à \struck{équivalence}
isomorphisme près), ou le cas échéant, que
\[
  \operatorname{Fonc\,exacts}(\mathcal{M}, \mathcal{C}) \longrightarrow
  \operatorname{Fonc\,add}(\mathcal{V}^{\circ}, \mathcal{C})
\]
soit une équivalence de catégories (ce qui

\page{144}
caractérise $\mathcal{M}$ à équivalence près). De plus on aimerait donner des
conditions moyennant lesquelles le foncteur $\underline{h}$ est
\emph{pleinement fidèle}.
\note{la « $\circ$ » en exposant de $\mathcal{V}$ (catégorie opposée) est une
lecture : il écrit $\mathcal{V}^{\circ}$ ou $\mathcal{V}^{*}$, et l'astérisque
de « additive${}^{*}$ » renvoie peut-être à la même marque ; « ${}^{2}$ »
souligné renvoie à une note qui n'est pas sur la page}

On va supposer que la structure \struck{pré}additive sur $\mathcal{V}$ peut
s'enrichir en une structure \emph{$\mathbf{Q}$-linéaire}
\marginal{i.e.\ en une structure d'anneau \emph{$k$-linéaire},
i.e.\ \uncertain{on a} $\mathbf{Q} \to \operatorname{End}\operatorname{id}_{\mathcal{V}}$}
[N.B.\ si $\mathbf{Q} = \mathbf{Q}$, \uncertain{cet enrichissement} est
évidemment unique et déterminé de façon unique]. On suppose que
\struck{\ill{} les \ill{} des \ill{}}
\struck{$\operatorname{Fl}(X,Y)$ \ill{} \ill{} finis sur $k$, et}
que pour tout $X$, $\operatorname{Fl}(X,X)$ est une algèbre \emph{semi-simple}
\uncertain{\emph{de dim finie}} sur $k$. \uncertain{Nous allons} alors construire
une catégorie $\mathcal{M}^{\#}$ de la façon suivante. Un objet de
$\mathcal{M}^{\#}$ est un couple $(X, \pi)$, où $X \in \operatorname{Ob}\mathcal{V}$,
$\pi \in \operatorname{Fl}(X,X)$, $\pi^{2} = \pi$ ($\pi$ est un \emph{projecteur}
de $X$). Un homomorphisme $(X,\pi) \to (X',\pi')$ est un homomorphisme
$f \colon X \to X'$ tel que $\pi' f = f \pi$, \emph{modulo} les $f$ tels que
$f\pi = \pi' f = 0$. La composition se fait de façon évidente. D'où une
catégorie $\mathcal{M}^{\#}$, manifestement additive, et un foncteur additif
\note{le premier $\mathbf{Q}$ du crochet est écrit gras, le second en lettre
barrée ; la lettre de l'enrichissement hésite entre $\mathbf{Q}$ et $k$, et la
marge dit « $k$-linéaire »}


\begin{tikzcd}[column sep=small, row sep=small]
  X \arrow[r, "\pi"] \arrow[d, "f"'] & X \arrow[d, "f"] \\
  X' \arrow[r, "\pi'"] \arrow[d, "g"'] & X' \arrow[d, "g"] \\
  X'' \arrow[r, "\pi''"] & X''
\end{tikzcd}

\note{diagramme dessiné dans la marge inférieure gauche, en face des lignes sur
la composition}

\page{145}
évident, d'ailleurs pleinement fidèle
\[
  \underline{h}^{\#} \colon \mathcal{V} \longrightarrow \mathcal{M}^{\#}
\]
en posant $\underline{h}^{\#}(X) = (X, \operatorname{id}_{X})$. Pour toute
catégorie abélienne $\mathcal{C}$, le foncteur
$\varphi \colon T \rightsquigarrow T \circ \underline{h}$
\[
  \operatorname{Fonc\,add}(\mathcal{M}, \mathcal{C})
  \xrightarrow{\ \varphi\ } \operatorname{Fonc\,add}(\mathcal{V}, \mathcal{C})
\]
est une \struck{isomorphisme} \struck{\ill{}} équivalence. Pour le voir, on
construit un foncteur quasi-inverse
\struck{1°) Pleinement fidèle. Un}
\marginal{il suffit que $\mathcal{C}$ soit une catégorie additive où les
projecteurs sont \uncertain{réalisables}}
\[
  \psi \colon \operatorname{Fonc\,add}(\mathcal{V}, \mathcal{C})
  \longrightarrow \operatorname{Fonc\,add}(\mathcal{M}, \mathcal{C})
\]
en associant à tout foncteur \uncertain{additif} $H \colon \mathcal{V} \to
\mathcal{C}$ son extension canonique par
\[
  \psi(H)(X,\pi) = \operatorname{Im} H(\pi)
\]
avec la loi fonctorielle évidente. On voit aussitôt que l'on a des isom.\ can.
\[
  \varphi\psi(H) \simeq H, \qquad \psi\varphi(T) \simeq T .
\]
\note{la marge gauche, écrite en oblique, accroche l'hypothèse « catégorie
abélienne » de la ligne ; ici $\underline{h}$ et $\operatorname{Fonc\,add}$ sont
écrits sans l'opposé de la page 143}

Pour l'instant cette construction et propriété universelle de $\mathcal{M}$
n'utilise pas les hypothèses particulières faites sur $\mathcal{V}$ ; elle
répond au Pb universel \ill{} (dans \uncertain{catégories additives données})
pour tous les projecteurs \uncertain{réalisables}.
\note{l'incise « (dans catégories additives données) » est ajoutée en
interligne ; un long trait, qui court sous « tous les projecteurs », la relie
peut-être à la marge oblique de la page}
Mais les hypothèses faites sur $\mathcal{V}$ vont impliquer que $\mathcal{M}$
est une catégorie \emph{abélienne} \emph{semi-simple} (i.e.\ tout objet est
\page{146}
\uncertain{provient} \ill{} \ill{} \uncertain{nombre} \ill{} : un objet $P$ de la
forme $\underline{h}(X)$, \struck{dans} $\operatorname{Fl}(P,P) \simeq
\operatorname{Fl}(X,X)$, \struck{dans \ill{} de \ill{} \ill{}}
\struck{$\underline{h}(\alpha)$, où $\alpha \in \operatorname{Fl}(X,X)$. Or
comme} $\operatorname{Fl}(P,P)$ est un \emph{anneau semi-simple} $U$,
\struck{On}
\struck{[Pour un tel anneau, dont idéal opérant \ill{}]}
\struck{[Alors $P$ est \ill{} \ill{} objet d'une catégorie \ill{} se décompose
suivant les composantes simples de $U = \operatorname{Fl}(P,P)$ \ill{} \ill{}
\ill{}]}
\note{deux passages successifs, encadrés et barrés, dont on ne lit que des
bribes ; « $\operatorname{Fl}(P,P) \simeq \operatorname{Fl}(X,X)$ » est ajouté
en interligne}

Or lorsque on a un anneau de ce type \uncertain{opérant} sur un objet d'une
catégorie abélienne, \uncertain{(\ill{})}, on \uncertain{voit}
\struck{\ill{} complexe}
que pour tout $u \in U$, \struck{il existe des idempotents} \emph{orthogonaux}
\emph{complémentaires} $\pi, \pi'$ de $U$,
\[
  (\pi\pi' = \pi'\pi = 0,\quad \pi + \pi' = 1,\quad \pi^{2} = \pi,\quad
  \pi'^{2} = \pi')
\]
\struck{tels que}
et $\varpi, \varpi'$ tels que $u\pi' = 0$, $\varpi' u = 0$, et $v \in U$ tel
que
$vu$ soit un \struck{projecteur} \uncertain{idempotent} $\pi$,
$uv$ soit $\varpi$, et on a les relations
$\pi v = v$, $u\pi = u$, $\varpi u = u$, $v\varpi = v$ ; de cette façon, on
gagne donc.
\note{les quatre relations sont écrites dans une bulle, plus bas sur la page,
rattachée par un trait à « et on a les relations »}


\begin{tikzcd}[column sep=small, row sep=small]
  P \simeq \pi P + \pi' P \arrow[r] & \pi P \arrow[d, "\wr"] \\
  Q \simeq \varpi Q + \varpi' Q & \varpi Q \arrow[l]
\end{tikzcd}

\[
\left\{
\begin{array}{l}
  vu = \text{projecteur } \pi \\
  uv = \text{projecteur } \varpi
\end{array}
\right.
\]


\begin{tikzcd}[column sep=small, row sep=small]
  P \arrow[rr, "u"] \arrow[d] & & Q \\
  P' \arrow[rr, "\sim"] & & Q' \arrow[u]
\end{tikzcd}

\note{dans la marge gauche, un premier dessin : $P \xrightarrow{u} Q$ avec une
flèche $v$ en retour tracée au-dessus, et $P \to P' \simeq Q' \to Q$ ; en
dessous, un second schéma plus confus où figurent $P' + P''$, $Q' + Q''$ et
une flèche épaisse $P' \to Q'$, qui ne se laisse pas redessiner sûrement}

Pour les conditions de chaînes, on utilise le fait que les
$\operatorname{Fl}(X,X)$ sont \struck{de degré} de dim.\ finie sur $k$. …

\page{147}
cette catégorie est semi-simple). Si donc $\mathcal{C}$ est une catégorie
abélienne, les foncteurs additifs de $\mathcal{M}$ dans $\mathcal{C}$ sont
automat.\ exacts. Cela prouve que le « Pb universel » \uncertain{(avec grain
de sel)} [en termes de catégories abéliennes (Variante 2)] est résolu par
$\mathcal{M}$.

Prouvons donc que $\mathcal{M}$ est une catégorie \emph{abélienne}
\emph{semi-simple}. Cela équivaut à ceci :

\struck{Tout monomorphisme est \ill{}}
\begin{enumerate}
  \item[(i)] Tout objet \struck{\ill{}} satisfait condition des chaînes
  ascendante et descendante pour les sous-objets
  \struck{i.e.\ l'ens.\ des \ill{}}
  \item[(ii)] Tout monomorphisme $u \colon P \to Q$ est « direct »,
  i.e.\ isom.\ à un monomorphisme \uncertain{simple}
  \[
    A \times B \longrightarrow B \times C, \qquad
    \operatorname{inj}_{1} \circ \operatorname{pr}_{2} .
  \]
\end{enumerate}
\note{la formule $\operatorname{inj}_{1} \circ \operatorname{pr}_{2}$ est écrite
seule sous la flèche ; telle quelle elle n'est pas un monomorphisme de
$A \times B$, et on attendrait l'injection du facteur commun $B$}

Pour vérifier (ii), \struck{\ill{}} considérons $R = P \times Q$,
et l'endomorphisme \struck{$\operatorname{pr}_{1}$}
$v = \operatorname{inj}_{2}\, u \operatorname{pr}_{2}$ de $R$. On voit
aisément que \struck{\ill{}} (ii) \uncertain{sont vraies} pour $u$
\uncertain{si et seulement si} \ill{} pour $v$, donc on
\uncertain{peut} supposer $P = Q$. D'autre part, soit $S$ un autre objet, et
soit $R = P + S$, $v = u \oplus \operatorname{id}_{S}$, alors on constate que
la propriété \uncertain{est vraie} pour $u$ si elle l'est pour $v$. Cela
\ill{},
\marginal{donnons-nous \ill{} \ill{} \ill{} \ill{} il est innocent de modifier
\ill{} \ill{}}
\note{la note marginale, écrite en oblique dans l'angle inférieur gauche,
accompagne un petit schéma où se lisent $P \to Q$, $P'$, $Q + Q'$ et
$P \oplus P' \to Q$}

\page{148}
On suppose donné trois catégories \uncertain{(additives)}
$\mathcal{V}, \mathcal{V}', \mathcal{V}''$, d'où $\mathcal{M}, \mathcal{M}',
\mathcal{M}''$, on suppose donné un foncteur biadditif
$\mathcal{V} \times \mathcal{V}' \to \mathcal{V}''$. Alors il se prolonge de
façon essent.\ unique en un foncteur biadditif
$\mathcal{M} \times \mathcal{M}' \to \mathcal{M}''$. En fait, on voit que la
propriété \uncertain{(première)} universelle qui définit $\mathcal{M}$ en
termes de $\mathcal{V}$, s'étend aux bifoncteurs, …

On utilisera cette construction surtout pour un bifoncteur
$\mathcal{V} \times \mathcal{V} \to \mathcal{V}$ « symétrique » et
« associatif » (et « unitaire »), noté \struck{$\otimes$} $\boxtimes$. Alors on
trouve un bifoncteur de même nature dans $\mathcal{M}$, noté $\otimes$. Si $k$
est de car.\ 0 (p.ex.\ $\mathbf{Q}$) on peut construire donc toutes les
opérations tensorielles \uncertain{(covariantes)} dans $\mathcal{M}$ [en
plusieurs arguments], par exemple on peut définir les
$\bigotimes^{n} M$, $\bigwedge^{n} M$, \struck{etc} $S^{n}(M)$ etc …
\note{le signe du bifoncteur sur $\mathcal{V}$ est un $\otimes$ surchargé puis
remplacé par un produit tensoriel dans un carré ; les exposants $n$ sur
$\bigotimes$ et $\bigwedge$ sont des lectures}
\page{149}
Considérons un foncteur additif

\begin{tikzcd}
  \mathcal{V} \arrow[rr, "H"] \arrow[dr, "\underline{h}"'] & & \mathcal{C} \\
  & \mathcal{M} \arrow[ur, "T"'] &
\end{tikzcd}

($\mathcal{C}$ cat.\ abélienne) ; comment exprimer que $T$ est \emph{fidèle} ?
Comme $T$ est exact, il revient au même de dire que
\struck{\ill{}}
\[
\left\{
\begin{array}{l}
  T \text{ fidèle} \\
  T(M) = 0 \Longrightarrow M = 0 \\
  T \text{ conservatif} \\
  T \colon \operatorname{End}(M,M) \to \operatorname{End}(T(M),T(M))
  \text{ injectif}
\end{array}
\right.
\]
\struck{Cela s'exprime aussi par la propriété}
Pour ceci il f.\ et s.\ que \uncertain{$M$} soit fidèle
\note{la phrase s'arrête là, soulignée, au-dessus d'un trait qui barre toute la
largeur du feuillet ; la lettre grasse « $M$ » est peut-être $\mathcal{M}$ ou
$H$, et « soit fidèle » semble inachevé}

\textbf{Exemple} $k$ corps alg.\ clos.
\[
\begin{array}{l}
  \mathcal{V}_{0} \quad \text{catégorie des schémas projectifs lisses sur } k \\
  \quad\downarrow \\
  \mathcal{V} \quad \text{catégorie ayant mêmes objets que } \mathcal{V}_{0},
  \text{ mais}
\end{array}
\]
les morphismes \uncertain{sont}
\[
  \operatorname{Fl}(X,Y) = \bigoplus_{i,j} C^{m_{j}}(X_{i} \times Y_{j})
\]
où les $X_{i}$ sont les comp.\ connexes de $X$, les $Y_{j}$ celles de $Y$,
$\dim Y_{j} = m_{j}$, \uncertain{où} $C^{q}(Z)$ désigne les classes de cycles,
\uncertain{à coeff.\ dans $\mathbf{Q}$}, \uncertain{modulo} \ill{}
$\mathbf{Q}_{\ell}$-homologie, de codim.\ $q$ sur $Z$ \ill{} $L$.
\uncertain{La} composition se faisant suivant
\note{le signe de somme, d'abord tracé autrement, est surchargé en
$\bigoplus$ ; « à coeff.\ dans $\mathbf{Q}$ » est écrit en interligne, sur un
mot biffé ; la dernière ligne, au ras du bas du feuillet, porte une lettre
barrée (peut-être $\mathbf{Z}$) et « $L$ », que nous ne savons pas lire}

\page{150}
la composition connue des classes de correspondance. On a catégorie additive
(objet nul est $\varnothing$, somme est $\amalg$).
Si on admet que l'équivalence homologique = équivalence numérique, alors on
déduit \uncertain{(par th.\ de finitude de la cohomologie)} que les
$\operatorname{Fl}(X,Y)$ sont des vectoriels de dim.\ finie sur $\mathbf{Q}$ ;
si on admet l'inégalité de Hodge, les $\operatorname{Fl}(X,X)$ deviennent une
algèbre à involution du type de Castelnuovo-Weil …\ donc est semi-simple.
D'où une catégorie \struck{de} $\mathcal{M}$, mais on ne considère cette
catégorie \uncertain{qu'en} $\mathcal{V}$ ; \uncertain{il} ne \uncertain{reste}
que \ill{} \ill{} un foncteur \emph{contravariant} additif universel
\[
  \mathcal{V}^{\circ} \xrightarrow{\ \underline{h}\ } \mathcal{M}
\]
D'ailleurs, on a un foncteur évident
\[
  H^{*}(\ ,\mathbf{Q}_{\ell}) \colon \mathcal{V}^{\circ} \longrightarrow
  \operatorname{Mod}_{\mathrm{t.f.}}(\mathbf{Q}_{\ell})
\]
en associant à tout $X \in \operatorname{Ob}\mathcal{V}$ le module
$H^{*}(X, \mathbf{Q}_{\ell})$. Ce foncteur est \emph{fidèle}, par définition
\uncertain{même} des \uncertain{Hom} dans $\mathcal{V}$ (+ dualité …).
Il \uncertain{définit} un foncteur exact fidèle \uncertain{conserv.}
\[
  T_{\ell} \colon \mathcal{M} \longrightarrow
  \operatorname{Mod}_{\mathrm{t.f.}}(\mathbf{Q}_{\ell})
\]
appelé \emph{foncteur de Tate}.
\note{« par th.\ de finitude de la cohomologie » est une lecture très incertaine
d'une ligne rapide ; le nom propre est écrit « Castelnuovo-Weil »}

\page{151}
Autres foncteurs \struck{canoniques} \uncertain{remarquables} sur
$\mathcal{M}$ : foncteurs de Hodge, de \struck{Hodge} De Rham (fidèles
en car.\ 0), de \emph{Betti} (si $k = \mathbf{C}$). Il y a d'ailleurs des structures
supplémentaires importantes à mettre sur ces foncteurs …

On a un bifoncteur $\mathcal{V} \times \mathcal{V} \to \mathcal{V}$
\[
  X, Y \rightsquigarrow X \times Y
\]
\struck{\ill{}} qui est associatif et commutatif, d'où une structure
multiplicative dans $\mathcal{M}$. La formule de Künneth en cohomologie
indique que $T_{\ell}$ est un foncteur qui commute \struck{aux} à
l'opération $\otimes$, donc aux opérations tensorielles.

\note{trois traits horizontaux séparent ce qui précède de ce qui suit}

Revenons à la situation abstraite générale. Nous supposons maintenant que
\uncertain{pour tout} $X \in \operatorname{Ob}\mathcal{V}$, on se donne des
\emph{projecteurs} $(\pi_{i} \in \operatorname{Fl}(X,X)$ $(i \in \mathbf{N}))$,
\uncertain{presque tous nuls},
(de somme $1$, deux à deux orthogonaux, fonctoriels en

\page{152}
$X$. En d'autres termes, on suppose donnés \struck{un \uncertain{homomorphisme}
\uncertain{d'anneaux} topologiques} \uncertain{éléments}
\[
  \pi_{i} \in \operatorname{End}\operatorname{id}_{\mathcal{V}}, \quad
  \text{avec } \pi_{i}^{2} = \pi_{i},\ \pi_{i}\pi_{j} = \pi_{j}\pi_{i} = 0
  \text{ si } i \neq j,
\]
\struck{$k[[\pi_{1}, \ldots, \pi_{n}, \ldots]] / (\pi_{i}^{2} - \pi_{i},
\pi_{i}\pi_{j}, \Sigma\pi_{i} - 1)$}
et $\Sigma\pi_{i} = 1$ (somme infinie !). Alors les $\pi_{i}$ s'étendent
à $\mathcal{M}$, qui devient une catégorie graduée, i.e.\ \uncertain{produit}
« somme » de catégories $\mathcal{M}^{i}$, dont les objets sont appelés
« motifs \uncertain{effectifs} de poids $i$ ».
\note{« effectifs » est ajouté en interligne ; la lettre « $\xi$ » qu'on devine
sous l'exposant de $\mathcal{M}^{i}$ ne se lit pas sûrement}

\uncertain{Si} on suppose de plus que \struck{\ill{}} pour $X \boxtimes Y$, on a
\[
  \pi_{n}(X \boxtimes Y) = \sum_{i+j=n} \pi_{i}(X) \boxtimes \pi_{j}(Y) ;
\]
ceci implique alors que $\mathcal{M}^{i} \times \mathcal{M}^{j} \to
\mathcal{M}^{i+j}$ par l'accouplement $\otimes$, et fournit donc un lien
étroit entre les $\mathcal{M}^{i}$. Ainsi, le groupe \uncertain{des classes} $K$ de $\mathcal{M}$ est un \emph{$\lambda$-anneau gradué}.
L'objet unité $\Lambda(0)$ est de degré $0$.
\struck{On va supposer que c'est le seul objet simple de ce degré, et que son
anneau \ill{} \ill{} \ill{} identique à $\mathbf{Q}$.}
\note{les deux lignes biffées se lisent en partie : « le seul objet simple »,
« de ce degré », « identique à $\mathbf{Q}$ »}

\struck{\ill{}} Je suppose maintenant \uncertain{fixé}, \struck{dans}
\uncertain{bien}, un objet \struck{\ill{}} $\Lambda(1)$ de $\mathcal{M}$
(\uncertain{de degré $p \geq 1$}), tel que, posant \struck{\ill{}}
$M(1) = M \otimes \Lambda(1)$, \struck{\ill{} \ill{}} le
\note{« de degré $p \geq 1$ » est une insertion en interligne, reliée par un
trait à la parenthèse ouvrante devant « tel » ; la phrase continue au feuillet
suivant}
\page{153}
foncteur $M \rightsquigarrow M(1)$ \emph{soit pleinement fidèle}. On itère
\uncertain{pour avoir} les $M \rightsquigarrow M(n) = M \otimes \Lambda(n)$,
$\Lambda(n) \simeq \Lambda(1)^{\otimes n}$. Notons que
\[
  \operatorname{Hom}(\Lambda(n), \Lambda(n)) \simeq
  \operatorname{Hom}(\Lambda(0), \Lambda(0)) \simeq \mathbf{Q}
\]
On introduit alors formellement aussi les $M(n)$, $n \in \mathbf{Z}$, en
considérant la catégorie \emph{pseudo-$\varinjlim$} des copies de
$\mathcal{M}$, par les foncteurs $M(1)$. Ainsi tout objet de la nouvelle
catégorie $\widetilde{\mathcal{M}}$ \struck{« motifs \ill{} »} s'écrit
\[
  M(i) \qquad i \in \mathbf{Z},\ M \in \operatorname{Ob}\mathcal{M},
\]
avec
\[
  \operatorname{Hom}(M(i), N(j)) = \operatorname{Hom}(M(i+r), N(j+r))
\]
où $r$ est pris tel que $i+r, j+r \geq 0$, de sorte que $M(i+r)$, $N(j+r)$ est
\uncertain{défini} \uncertain{comme} \uncertain{motif}. La catégorie
$\widetilde{\mathcal{M}}$ est encore graduée (de \uncertain{manière} évidente),
\uncertain{et} périodique de période $p$ par les foncteurs
$M \rightsquigarrow M(1)$. En \uncertain{dim.}\ $i \geq 0$ (\uncertain{ou
quelconque}, \ill{} \ill{} \uncertain{négatif}) \uncertain{la} \ill{} s'écrit
comme
\[
  \operatorname{Ps}\varinjlim_{k} \mathcal{M}^{i+kp}, \qquad
  \text{où } \mathcal{M}^{i+kp} \longrightarrow \mathcal{M}^{i+(k+1)p}
  \text{ est le foncteur pl.\ fid.\ } M \rightsquigarrow M(1).
\]
\marginal{$\Lambda(1)$ de \uncertain{poids} $p$}
\marginal{N.B.\ \uncertain{Un} \uncertain{sous}-motif d'un motif effectif
est effectif.}
\note{la note « $\Lambda(1)$ de poids $p$ » est dans la marge gauche, en face
de la ligne sur $M(i+r)$ ; le N.B.\ est écrit en petit dans l'angle inférieur
droit. « $\operatorname{Ps}\varinjlim$ » abrège la pseudo-limite inductive
nommée plus haut}

\page{154}
\struck{Décomposition canonique}

\textbf{Niveau d'un motif}
\note{titre souligné de sa main, sous un premier titre biffé qu'on lit
« Décomposition canonique »}

a) \emph{Motif effectif.} $M$ de degré $i \geq 0$ est dit de niveau
$\leq i - kp$ s'il est isomorphe à un \struck{\ill{}} motif de la forme
$M'(k)$, où $M'$ est un motif effectif de degré $i - kp$.
Les motifs effectifs, de degré $i \geq 0$ fixé, de niveau $\leq i - kp$, forment
une sous-catégorie \emph{épaisse} de $\mathcal{M}^{i}$, donc tout objet $M$ de
$\mathcal{M}^{i}$ se décompose canoniquement en une somme d'un
\uncertain{sous-objet} maximal \uncertain{(\ill{})}
$M'$ de niveau $\leq i - kp$, et d'un \uncertain{autre} \uncertain{sous-objet}
\ill{}.
\note{« $M$ » est écrit en exposant après « tout objet » ; « maximal » est
ajouté en interligne au-dessus de « objet », et une insertion reliée à la
dernière parenthèse, peut-être « autre sous-objet », ne se lit pas
sûrement}

Ainsi,
\[
  \mathcal{M}^{i} \simeq \mathcal{P}^{i} \times \mathcal{P}^{i-p}(1) \times
  \mathcal{P}^{i-2p}(2) \times \cdots ,
\]
où $\mathcal{P}^{j}$ désigne la catégorie des motifs effectifs de degré $j$
\struck{\ill{}} dont aucun sous-motif $\neq 0$ n'est de niveau $\leq j - p$.
\marginal{projections $K_{i,k}$}
\marginal{motifs \uncertain{primitifs} de niveau $j$}
\note{les deux notes sont dans la marge gauche ; la seconde, écrite en oblique,
est soulignée ; la lettre ronde nommant ces catégories est lue
$\mathcal{P}$ (pour « primitifs » ?), lecture incertaine. La note
« projections $K_{i,k}$ » se rapporte à la décomposition en produit}

Cela donne donc une décomposition de $\mathcal{M}^{*} = \mathcal{M}$ en termes
des $\mathcal{P}^{i}$ (pour les structures multiplicative et additive). La
structure \uncertain{multiplicative} \uncertain{se} \uncertain{traduit} \ill{}
\uncertain{par} les foncteurs
\[
  \mathcal{Q}^{i,j,k} \colon \mathcal{P}^{i} \times \mathcal{P}^{j}
  \longrightarrow \mathcal{P}^{i+j-kp}
  \qquad
  \mathcal{Q}^{i,j,k}(M,N) = K_{i,k}(M \otimes N)(k) .
\]
\note{la lettre du foncteur, ronde et surchargée, est lue $\mathcal{Q}$ sans
certitude ; dans la formule de droite, un premier essai biffé précède
« $K_{i,k}(M \otimes N)(k)$ »}

b) \emph{Niveau d'un motif quelconque} (pas nécessairement effectif). On le
définit en se ramenant au cas précédent …

\page{155}
\textbf{Duals et \uncertain{Homs} de motifs}

On supposera que (\uncertain{$p = 2$ et que}) pour tout motif effectif $M$, de
degré $i \geq 0$, $\underline{\operatorname{Hom}}(M, \Lambda(-i))$ existe
(i.e.\ \uncertain{est} \uncertain{un} motif eff.),
\note{« $p = 2$ et que » est ajouté en interligne au-dessus de « pour tout » ;
l'insertion « est un motif eff. » court au-dessus de la fin de la ligne
suivante et se rattache à la parenthèse « (i.e. »}
où $\underline{\operatorname{Hom}}(M,N)$ est défini axiomatiquement par
\[
  \operatorname{Hom}(P, \underline{\operatorname{Hom}}(M,N)) \simeq
  \operatorname{Hom}(P \otimes M, N)
\]
\note{le $\operatorname{Hom}$ interne s'écrit d'une majuscule cursive grasse,
distincte du $\operatorname{Hom}$ ordinaire à droite ; nous le rendons
$\underline{\operatorname{Hom}}$}
On constate que si $M$, $N$ sont des motifs effectifs, de degrés $i$, $j$,
tels que $\underline{\operatorname{Hom}}(M,N)$ existe, alors
$\underline{\operatorname{Hom}}(M,N)$ est de degré $j - i$. Ici, posant
\[
  M^{*} = \underline{\operatorname{Hom}}(M, \Lambda(-i)) \qquad
  \deg M = i
\]
on trouve que $\deg M^{*} = \deg M$. On trouve d'ailleurs que si
$\underline{\operatorname{Hom}}(M,N)$ existe dans $\mathcal{M}$, il existe
aussi dans $\widehat{\mathcal{M}}$ et est le même.
\note{la lettre porte ici un accent circonflexe, $\widehat{\mathcal{M}}$, et non
le tilde de la catégorie $\widetilde{\mathcal{M}}$ de la page 153 ; nous
gardons ce qui est écrit}

On voit que pour tout $M$ de degré $i$ et tout $N$,
$\underline{\operatorname{Hom}}(M, N(-i))$ existe, car on a
\[
  \underline{\operatorname{Hom}}(M, N(-i)) \simeq M^{*} \otimes N
\]
(vérification facile). \struck{D'où} Cela prouve que pour tout $N$ (\uncertain{non}
\uncertain{nécess.} \uncertain{effectif})

\page{156}
$\underline{\operatorname{Hom}}(M, N(j))$ existe pour $j$ grand, dans la
catégorie des motifs effectifs. Enfin, $\underline{\operatorname{Hom}}(M,N)$
existe dans la catégorie de tous les motifs, comme on voit en prenant
$\underline{\operatorname{Hom}}(M, N(j))(-j)$. On arrive ainsi à construire
$\underline{\operatorname{Hom}}(M,N)$ pour $M, N \in
\operatorname{Ob}\widehat{\mathcal{M}}$ quelconques, avec les règles de calcul
évidentes. On définit aussi
$\check{M} = \underline{\operatorname{Hom}}(M, \Lambda(0))$, et on trouve
$\underline{\operatorname{Hom}}(M,N) \simeq \check{M} \otimes N$. Mais on fera
attention que si $M$ est effectif, $\check{M}$ ne l'est pas si
$\deg M \neq 0$.
\note{le dernier signe, « $\deg M = 0$ » ou « $\neq 0$ », n'est pas net ; le
sens exige $\neq$}

\[
  B(M,N) \simeq \operatorname{Hom}(M \otimes N, \Lambda(-i)) \simeq
  \operatorname{Hom}(M, N^{*}) \simeq \operatorname{Hom}(N, M^{*})
\]
« formes bilinéaires », avec $M$, $N$ de degré $i$.
Si $M = N$, on a la notion de forme alternée, forme symétrique.
\textbf{N.B.} Forme non dégénérée en $M$, i.e.\ $M \to N^{*}$ \uncertain{mono.}
$\Longleftrightarrow$ $N \to M^{*}$ est \struck{surjectif} \uncertain{épi.}
On définit de même forme non dégénérée en $N$.
\marginal{?}
Si $M = N$, les deux notions sont équivalentes \uncertain{car} \struck{\ill{}}
\uncertain{pour} \uncertain{un} \uncertain{homom.} $M \to M^{*}$,
\uncertain{dire} \uncertain{que} c'est \uncertain{mono.} \uncertain{ou}
\uncertain{épi.} \ill{} \uncertain{revient} \uncertain{au} \uncertain{même} \ill{}
\uncertain{(\ill{})}.
\note{un grand point d'interrogation, dans la marge gauche, porte sur ces deux
dernières lignes, qui sont d'une écriture très rapide ; la lettre $B$ de la
forme bilinéaire est un B cursif ouvert}

\page{157}
\textbf{Involutions.} Supposons $\alpha \in B(M,M)$ \struck{forme
\uncertain{fondamentale}} forme non dégénérée, symétrique ou \struck{\ill{}}
alternée, d'où \struck{Alors pour \ill{}} un \uncertain{isom.}
$\widetilde{\alpha} \colon M \xrightarrow{\sim} M^{*}$. Soit de même
$\beta \in B(N,N)$ d'où $\widetilde{\beta} \colon N \to N^{*}$
($M$, $N$ de deg.\ $i$). Soit $u \colon M \to N$, on définit alors
$u' \colon N \to M$ par \uncertain{commutativité} du diagramme ci-contre.

\begin{tikzcd}
  M \arrow[r, leftrightarrow, "u^{\sharp}"] \arrow[d, "\widetilde{\alpha}"'] & N \arrow[d, "\widetilde{\beta}"] \\
  M^{*} & N^{*} \arrow[l, "{}^{t}u^{*}"']
\end{tikzcd}

\note{la flèche du haut porte une pointe à chaque bout, avec une seule
étiquette qu'on lit $u^{\sharp}$ (pour les deux sens $u$ et $u'$ ?) ; sous la
flèche du bas, un trait mène au mot « transposée », lecture douteuse ; la
définition de $u'$ ne distingue pas ici $u'$ de $u^{\sharp}$}

On a les relations habituelles
\[
\begin{aligned}
  (u \pm v)' &= u' \pm v' \\
  (uv)' &= v'u' \\
  1' &= 1 \\
  (u')' &= u
\end{aligned}
\]
si $\alpha$ et $\beta$ sont tous deux \struck{\uncertain{symétriques}} ou
tous deux alternés.
\note{la condition sur $\alpha$, $\beta$ est écrite à droite de la dernière
relation, en deux lignes serrées ; le mot biffé, « symétriques » peut-être,
est remplacé en dessous par « ou tous deux alternés »}

En particulier sur $\operatorname{End}(M)$ \uncertain{on} a une structure
d'algèbre involutive sur $\mathbf{Q}$. \struck{\ill{}}
Si $\alpha$ est remplacé par $\alpha_{1}$, $\beta$ par $\beta_{1}$, avec
\[
  \widetilde{\alpha}_{1} = \widetilde{\alpha} \circ A, \qquad
  \widetilde{\beta}_{1} = \widetilde{\beta} \circ B \qquad
  (A \in \operatorname{End}(M),\ B \in \operatorname{End}(N))
\]
\uncertain{alors} la nouvelle « adjonction » \uncertain{sera}
\[
  u'' = \widetilde{\alpha}_{1}^{-1} \circ u^{*} \circ \widetilde{\beta}_{1}
  = A^{-1}\widetilde{\alpha}^{-1} u^{*} \widetilde{\beta} B \quad \text{i.e.}
\]
\[
  u'' = A^{-1} u' B
\]
En particulier, si $M = N$, on a \uncertain{pour} $A$ \uncertain{une}
\uncertain{nouvelle} involution définie \uncertain{modulo}
\struck{\uncertain{automorphismes} intérieurs} \uncertain{modification de
nature \ill{}} : $u \mapsto A^{-1} u' A$, où $A$ est un \uncertain{élément}
autoadjoint (?).
\note{dans la première formule, un premier symbole biffé précède
$\widetilde{\alpha}_{1}$ ; le transposé y est écrit $u^{*}$ et non
${}^{t}u^{*}$ comme sur le diagramme. La fin de page, après « modulo », est
surchargée : « automorphismes intérieurs » est biffé et remplacé en
interligne par deux ou trois mots peu lisibles}

\page{158}
\textbf{Formes définies positives sur les motifs}
\note{titre souligné de sa main, en tête du feuillet, d'une encre plus noire}

\uncertain{Lui} indiquer déjà que les $\operatorname{End}(h^{i}(X))$ sont des
\emph{algèbres de Riemann}, et par suite semi-simples, (ce qui permet de
construire la catégorie des motifs comme catégorie \emph{abélienne}).
\emph{De plus}, cette catégorie des motifs $\mathcal{M}$, et sa complétion
$\overline{\mathcal{M}}$, acquièrent d'importantes structures
\uncertain{supplémentaires}.
\note{« ½ simples » est son abréviation de « semi-simples » ; le premier mot,
lu « Lui », suggère une note à l'usage d'un correspondant, sans que la page
le nomme ; la complétion est ici surlignée, $\overline{\mathcal{M}}$, et non
chapeautée comme aux pages 155-156}

Discutons d'abord le degré d'indétermination possible du système
$\varphi_{X}$. Pour ceci, notons que \struck{\ill{}} \uncertain{si} on
\uncertain{remplace} les $(\varphi_{i})$ par $(\varepsilon_{i}\varphi_{i})$
($\varepsilon_{i}$ \emph{signes} \uncertain{dépendant} de $i$). Les${}^{*}$
axiomes \ill{} \ill{} : (i) et (iii) \uncertain{restent} vérifiés, mais (ii) seulement si on a
$\varepsilon_{i}\varepsilon_{j} = \varepsilon_{i+j}$, donc si $\varepsilon_{i}$
\uncertain{dépend} \uncertain{de} $i$ sous \struck{la forme \ill{} $(\ldots)$}
$+1$ pour tout $i$, ou $\varepsilon^{i}$ \uncertain{--} il est
\uncertain{nécessaire} \uncertain{de} \uncertain{demander} \uncertain{ce}
\uncertain{qu'il} \uncertain{doit} \uncertain{être} \uncertain{en} degré $1$.
On imposera, dans le cas de la géom., qu'il corresponde aux formes
\uncertain{de} polarisation sur les V.A.
\note{les axiomes (i), (ii), (iii) du système $\varphi_{X}$ sont ceux que
pose la page 160 (voir plus loin) : le feuillet 160 précède logiquement
celui-ci ; l'astérisque de « Les${}^{*}$ »
n'a pas d'appel visible. Le raisonnement lisible : les changements de signe
$\varphi_{i} \mapsto \varepsilon_{i}\varphi_{i}$ respectant (ii) sont ceux où
$\varepsilon_{i}\varepsilon_{j} = \varepsilon_{i+j}$, i.e.\ $\varepsilon_{i} = 1$
pour tout $i$ ou $\varepsilon_{i} = \varepsilon^{i}$, et il reste à fixer le
signe en degré $1$ par les polarisations des variétés abéliennes. La fin de la
page est d'une écriture très rapide, et l'enchaînement proposé ici n'est pas
sûr}

\page{160}
Supposons que pour tout $X \in \mathcal{V}$ (on écrit \uncertain{que}
$\underline{h}(X) = (\underline{h}^{i}(X))$) \struck{soit \uncertain{donné}
\ill{}} un ensemble $\Phi_{X}$ de formes fondamentales $\{\varphi\}$,
i.e.\ que $\varphi = (\varphi_{i})$,
$\varphi_{i} \in \operatorname{Hom}(\underline{h}^{i}(X) \otimes
\underline{h}^{i}(X), \Lambda(-i))$, de façon que les conditions suivantes
soient vérifiées :
\note{la parenthèse « on écrit que $\underline{h}(X) = (\underline{h}^{i}(X))$ »
est une insertion, reliée par un trait à « $X \in \mathcal{V}$ » ; la lettre
de l'ensemble de formes, un $\mathcal{L}$ ou un $\Phi$ cursif, est lue
$\Phi$}

\begin{enumerate}
  \item[(i)] Si $\varphi \in \Phi_{X}$, $\psi \in \Phi_{Y}$, alors
  $\varphi \oplus \psi \in \Phi_{X \amalg Y}$.
  \item[(ii)] Si $\varphi \in \Phi_{X}$, $\psi \in \Phi_{Y}$,
  \struck{et $u \colon X \to Y$, d'où $\underline{h}(u) \colon
  \underline{h}^{i}(Y) \to \underline{h}^{i}(X)$, \uncertain{considérons}}
  alors $\varphi \otimes \psi \in \Phi_{X \boxtimes Y}$.
  \item[(iii)] Sur \struck{\ill{}} $h^{0}(X)$ \uncertain{(\uncertain{extension}
  de $\mathbf{Q}$)} \struck{\ill{} $\varphi(1)$ \ill{} $X = e$, donc}
  \struck{$\underline{h}^{0}(e) \simeq \underline{h}^{0}(e) \simeq \mathbf{Q}$,
  alors} \uncertain{fini} sur $\mathbf{Q}$, les $\varphi \in \Phi_{X}$ sont
  des formes \emph{définies positives} [cela suffit \struck{\ill{}} \uncertain{pour}
  $X = e$] \uncertain{le} \uncertain{supposer} pour \ill{} \uncertain{sans}
  \uncertain{doute} \uncertain{dans} \ill{} \ill{}
  \item[(iii)] Si $\varphi \in \Phi_{X}$, $\psi \in \Phi_{Y}$,
  $u \colon X \to Y$, d'où $u^{(i)} \colon h^{i}(Y) \to h^{i}(Y)$, prenant
  $s_{u}^{(i)} \colon h^{i}(Y) \to h^{i}(X)$ [\uncertain{relatif} à
  $\varphi$, $\psi$], on a
  \[
    \boxed{\operatorname{Tr}\bigl(s_{u}^{(i)} u^{(i)}\bigr) \geq 0,
    \ \text{égalité ssi } u^{(i)} = 0}
  \]
\end{enumerate}
\note{la condition (iii) est écrite deux fois : une première, dont le numéro
porte un signe qu'on ne sait lire (« (iii)° » ?), sur la positivité en degré
$0$, plusieurs fois surchargée ; puis un « (iii) » biffé, et un nouveau (iii),
la positivité de la trace, que la dernière formule encadre. Telle qu'écrite,
$u^{(i)}$ va de $h^{i}(Y)$ dans $h^{i}(Y)$ ; on attendrait
$h^{i}(Y) \to h^{i}(X)$ pour $\underline{h}(u)$ et $s_{u}^{(i)}$ en sens
inverse — nous gardons la page. $s_{u}^{(i)}$ est sans doute l'adjoint de
$u^{(i)}$ relativement à $\varphi$, $\psi$, au sens de la page 157}

\end{document}
